\scanpage{59}
% archive foliation: 58
% Sheet read after rotation through 180 degrees.

9.

4.3. - \uline{La politique de recherche des établissements}

Un établissement d'enseignement supérieur possède en général un Conseil scientifique. Il lui revient de définir les grandes lignes de la politique scientifique de l'établissement. L'objectif est clair :

\begin{itemize}
\item[-] compte tenu des moyens en bâtiments, personnel enseignant et technicien, comment utiliser les crédits spécifiques pour la recherche de façon optimum, pour faire la meilleure recherche, \unsure{celle} qui assure le renon de l'établissement et lui permet de s'associer avec les grands établissements de recherche, en premier lieu avec le C.N.R.S.
\item[-] Il appartient au Conseil scientifique de conduire la politique de recherche de l'établissement, sous la responsabilité du Président. La recherche est l'affaire de l'ETABLISSEMENT, chacun étant solidaire pour la mener à bien.
\end{itemize}

Une politique recherche est, typiquement, celle que fait naitre des équipes qui deviennent peu à peu des équipes associées avec le C.N.R.S. et avec l'I.N.S.E.R.M.. La collaboration prolongée de l'établissement à frais partagés avec un organisme de recherche est le critère de réussite d'une politique de recherche.

5 - \uline{La politique ``recherche'' de la D.E.S.U.R.}

5.1. \uline{La D.E.S.U.R. doit remplir une mission de complexe dans le cadre du plan-recherche lui permettant} :

a) d'agir ``en aval'' au niveau des établissements en respectant leur autonomie ; elle ne peut donner des ordres ni définir des programmes \uline{a priori} ; elle doit par contre, effectuer un contrôle a postériori, pour avoir des éléments de jugement sur l'activité de recherche des établissements.

b) d'agir ``en amont'' au niveau national en tenant sa place, parmi une dizaine d'organismes fortement hiérarchisés, pour obtenir sa part de l'enveloppe-recherche. La D.E.S.U.R. doit être ``crédible'' par la D.G.R.S.T.

